% !TEX program = xelatex
\documentclass[aspectratio=169,11pt]{beamer}

\usepackage{amsmath,amssymb,booktabs,tabularx}
\usepackage{CUHKBeamer}

\title[CUHK Beamer Template]{Research Presentation Title}
\subtitle{A concise subtitle for the project or paper}
\author[Your Name]{Your Name}
\institute[CUHK]{Department or Programme\\The Chinese University of Hong Kong}
\date{September 2026}

\begin{document}

\begin{frame}[plain]
  \titlepage
\end{frame}

\begin{frame}{Outline}
  \tableofcontents
\end{frame}

\section{Motivation}

\begin{frame}{A Clear Research Motivation}
  Strong presentations establish the problem before introducing the method. State the
  practical or scientific setting, isolate the unresolved question, and explain why it
  matters to the audience.

  \vspace{0.45cm}
  \begin{columns}[T,onlytextwidth]
    \begin{column}{.48\textwidth}
      \begin{block}{Context}
        Summarise the established knowledge and the assumptions your work inherits.
      \end{block}
    \end{column}
    \begin{column}{.48\textwidth}
      \begin{alertblock}{Research gap}
        Identify one precise limitation that the study is designed to address.
      \end{alertblock}
    \end{column}
  \end{columns}

  \vspace{0.25cm}
  \cuhkkeypoint{Express the research question in one sentence that can be tested.}
\end{frame}

\begin{frame}{Objectives and Contributions}
  This template supports short talks, seminars, proposals, and thesis defences. Replace
  the example statements below with claims supported by your own evidence.

  \vspace{0.35cm}
  \begin{enumerate}
    \item Define the problem and evaluation setting precisely.
    \item Present a method whose components can be inspected independently.
    \item Report results together with limitations and validation status.
  \end{enumerate}

  \vspace{0.25cm}
  \begin{exampleblock}{Recommended slide rhythm}
    Alternate text, diagrams, equations, and tables instead of repeating one layout.
  \end{exampleblock}
\end{frame}

\section{Method}

\begin{frame}{Method Overview}
  \begin{columns}[c,onlytextwidth]
    \begin{column}{.57\textwidth}
      \centering
      \begin{tikzpicture}[
        node distance=.38cm,
        box/.style={draw=CUHKPurple,very thick,rounded corners=2pt,
          fill=CUHKSoftGray,minimum width=1.65cm,minimum height=.75cm,align=center},
        arrow/.style={-{Latex[length=2.2mm]},thick,draw=CUHKGold}]
        \node[box] (input) {Input};
        \node[box,right=of input] (model) {Model};
        \node[box,right=of model] (output) {Output};
        \draw[arrow] (input) -- (model);
        \draw[arrow] (model) -- (output);
      \end{tikzpicture}
    \end{column}
    \begin{column}{.39\textwidth}
      \textbf{Input}\,---,Define the data and units.\\[6pt]
      \textbf{Model}\,---,Name the essential transformation.\\[6pt]
      \textbf{Output}\,---,Connect predictions to evaluation.
    \end{column}
  \end{columns}

  \vspace{0.55cm}
  \cuhkkeypoint{A diagram should expose the data flow, not decorate the slide.}
\end{frame}

\begin{frame}{Core Formulation}
  Let $x \in \mathbb{R}^{d}$ denote an observation and $y$ its target. A generic objective is
  \begin{equation*}
    \theta^{\star}=\arg\min_{\theta}\frac{1}{n}\sum_{i=1}^{n}
    \mathcal{L}\!\left(f_{\theta}(x_i),y_i\right)+\lambda\,\Omega(\theta).
  \end{equation*}

  \begin{columns}[T,onlytextwidth]
    \begin{column}{.48\textwidth}
      \begin{block}{Data term}
        $\mathcal{L}$ measures disagreement between the prediction and observation.
      \end{block}
    \end{column}
    \begin{column}{.48\textwidth}
      \begin{block}{Regularisation}
        $\Omega$ encodes an explicit preference; $\lambda$ controls its influence.
      \end{block}
    \end{column}
  \end{columns}
\end{frame}

\section{Results}

\begin{frame}{Results at a Glance}
  \centering
  \small
  \begin{tabularx}{.88\textwidth}{l *{3}{>{\centering\arraybackslash}X}}
    \toprule
    Method & Accuracy $\uparrow$ & Error $\downarrow$ & Runtime $\downarrow$ \\
    \midrule
    Baseline A & 78.4 & 0.214 & 42 min \\
    Baseline B & 81.7 & 0.186 & 55 min \\
    \textbf{Proposed} & \textbf{85.2} & \textbf{0.151} & \textbf{38 min} \\
    \bottomrule
  \end{tabularx}

  \vspace{0.55cm}
  \begin{alertblock}{Interpretation}
    Replace illustrative values with verified results and state the evaluation protocol.
  \end{alertblock}
\end{frame}

\begin{frame}{Conclusion}
  The final slide should return to the original research question and distinguish evidence
  from interpretation.

  \vspace{0.4cm}
  \begin{enumerate}
    \item State the central empirical finding.
    \item Explain its significance without overstating generality.
    \item Name the most important limitation or next experiment.
  \end{enumerate}

  \vfill
  \cuhkkeypoint{Thank you. Questions and discussion are welcome.}
\end{frame}

\end{document}
